\subsubsection{Schritt 1: Gaussian Blur}

Um einen mathematisch optimalen Kantendetektor herzuleiten, trifft Canny
eine Annahme über das Rauschen im Bild: \textit{,,We will assume that the
image consists of the edge and additive white Gaussian noise``}
\cite[S.~680]{canny1986}. Da dieses Rauschen fälschlicherweise als Kante
detektiert werden kann, muss das Bild zunächst geglättet werden. Canny
zeigt, dass der optimale Filter für dieses Rauschmodell durch die erste
Ableitung einer Gaußfunktion approximiert werden kann
\cite[S.~687]{canny1986}:

\begin{equation}
    \label{eq:gaussian}
    G(x) = \exp\left(-\frac{x^2}{2\sigma^2}\right)
\end{equation}

Der Parameter $\sigma$ kontrolliert die Stärke der Glättung. Zwischen
Rauschunterdrückung und Lokalisierungsgenauigkeit besteht jedoch ein
grundlegender Zielkonflikt, den Canny als
\textit{,,uncertainty principle``} zwischen Detektion und Lokalisierung
bezeichnet \cite[S.~684]{canny1986}.